\begin{remark}[5.11.3]
\label{IV.5.11.3}
The conclusion that $A^{(1)}$ is a finite $A$-algebra need no
longer hold if, in hypothesis {\rm(ii)} of
\sref{IV.5.11.2}, one assumes only that $A$ is catenary.

Indeed, consider the catenary local ring $A$ constructed in
\sref{IV.5.6.11}; its integral closure $A'$ (denoted there by
$B$) is a finite $A$-algebra.  If $A^{(1)}$ were also finite over
$A$, then, being contained in the fraction field of $A$, it would
be contained in $A'$.  On the other hand, in the notation of
\sref{IV.5.6.11}, every height-one prime ideal of $A$ is of the
form $\mathfrak pA$, where $\mathfrak p$ is a height-one prime of
$C$, and
\[
  A_{\mathfrak pA}=C_{\mathfrak p}.
\]
Moreover, $C_{\mathfrak p}=E_{\mathfrak p'}$, where
$\mathfrak p'$ is the unique prime of $E$ above $\mathfrak p$.
Thus $A_{\mathfrak pA}$ is integrally closed and consequently
contains $A'$.  By definition this gives
$A'\subset A^{(1)}$, so finiteness of $A^{(1)}$ would force
$A^{(1)}=A'$.  This is impossible: one of the two prime ideals of
$A'$ above the maximal ideal $\mathfrak nA$ of $A$ has height
$1$, whereas $\mathfrak nA$ has height $2$, contradicting
\sref{IV.5.10.17}[iv].
\end{remark}

\begin{corollary}[5.11.4]
\label{IV.5.11.4}
Let $X$ be a locally Noetherian prescheme, let $Z$ be a closed
subset of $X$, put $U=X-Z$, and let $i:U\to X$ be the canonical
injection.  Let $\sh F$ be a coherent $\sh O_U$-module.  A
necessary condition for $i_*(\sh F)$ to be a coherent
$\sh O_X$-module is that
\[
  \codim\bigl(\overline{\{x\}}\cap Z,\overline{\{x\}}\bigr)
  \geq2
  \qquad
  \text{for every }x\in\operatorname{Ass}(\sh F).
\]
This condition is sufficient in either of the following cases:
\begin{enumerate}[{\rm(i)}]
  \item $X$ is locally immersible in a regular scheme.
  \item $X$ is universally catenary and universally Japanese;
  by definition, the latter means that every point $x\in X$ has
  an affine open neighbourhood whose ring is universally
  Japanese \sref[0]{0.23.1.1}.
\end{enumerate}
\end{corollary}

\begin{proof}
By \sref[I]{I.9.4.7}, there is a coherent
$\sh O_X$-submodule $\sh G$ of $i_*(\sh F)$ such that
$\sh G|U=\sh F$.  Plainly,
\[
  \operatorname{Ass}(\sh G)
  \subset\operatorname{Ass}(\sh F)
  \subset\operatorname{Ass}\bigl(i_*(\sh F)\bigr),
\]
and
\[
  \operatorname{Ass}\bigl(i_*(\sh F)\bigr)
  =\operatorname{Ass}(\sh F)
\]
by \sref{IV.3.1.13}.  Hence
$\operatorname{Ass}(\sh G)=\operatorname{Ass}(\sh F)$.
It remains to apply \sref{IV.5.11.1} to $\sh G$, observing that
\[
  i_*(\sh F)=\sh H^0_{X/Z}(\sh G),
\]
and that, when $X$ is universally catenary and universally
Japanese, hypothesis {\rm(ii)} of \sref{IV.5.11.1} holds by
definition.
\end{proof}

\begin{corollary}[5.11.5]
\label{IV.5.11.5}
Let $X$ be a locally Noetherian prescheme, let $Z$ be a subset
of $X$ stable under specialization, and let $\sh F$ be a coherent
$\sh O_X$-module such that
$\operatorname{Ass}(\sh F)\subset X-Z$.  Consider the conditions:
\begin{enumerate}[{\rm a)}]
  \item $\sh H^0_{X/Z}(\sh F)$ is a coherent
  $\sh O_X$-module.
  \item[{\rm d)}] For every subset $T$ of $Z$ closed in $X$,
  $i_*(\sh F|X-T)$ is coherent, where
  $i:X-T\to X$ is the canonical injection.  It is enough here
  to consider an increasing filtered family of such closed
  subsets whose union is $Z$.
\end{enumerate}
Then {\rm a)} implies {\rm d)}.  If $X$ satisfies either
hypothesis {\rm(i)} or {\rm(ii)} of \sref{IV.5.11.1}, the two
conditions are equivalent.
\end{corollary}

\begin{proof}
The hypothesis and \sref{IV.3.1.13} give
\[
  \operatorname{Ass}\bigl(i_*(\sh F|X-T)\bigr)
  =\operatorname{Ass}(\sh F).
\]
It follows from \sref{IV.5.10.2} that the canonical maps
\[
  i_*(\sh F|X-T)\longrightarrow\sh H^0_{X/Z}(\sh F)
\]
are injective.  Condition {\rm a)} therefore implies {\rm d)}.

\oldpage[IV-2]{125}

Conversely, condition {\rm d)} and \sref{IV.5.11.1} imply, in
the notation of that proposition, that
\[
  \codim(T\cap Y_\alpha,Y_\alpha)\geq2.
\]
Since $Z$ is the union of its subsets closed in $X$, it follows
that
\[
  \codim(Z\cap Y_\alpha,Y_\alpha)\geq2.
\]
The last assertion now follows from \sref{IV.5.11.1}.
\end{proof}

\begin{corollary}[5.11.6]
\label{IV.5.11.6}
Let $A$ be a Noetherian ring and put $X=\Spec(A)$.  Consider the
following properties:
\begin{enumerate}[{\rm a)}]
  \item For every quotient domain $B$ of $A$, the ring $B^{(1)}$
  \sref{IV.5.11.2} is a finite $B$-algebra.
  \item For every coherent $\sh O_X$-module $\sh F$ and every
  subset $Z$ of $X$ stable under specialization such that
  \[
    \codim\bigl(\overline{\{x\}}\cap Z,\overline{\{x\}}\bigr)
    \geq2
  \]
  for every
  $x\in\operatorname{Ass}(\sh F)\cap(X-Z)$, the module
  $\sh H^0_{X/Z}(\sh F)$ is coherent.
  \item For every closed subset $T$ of $X$ and every coherent
  $\sh O_U$-module $\sh G$, where $U=X-T$, such that
  \[
    \codim\bigl(\overline{\{x\}}\cap T,\overline{\{x\}}\bigr)
    \geq2
    \qquad
    (x\in\operatorname{Ass}(\sh G)),
  \]
  the module $i_*(\sh G)$ is coherent, where $i:U\to X$ is the
  canonical injection.
  \item For every quotient domain $B$ of $A$ and every ideal
  $\mathfrak J$ of $B$ of height at least $2$, the ring
  \[
    \bigcap_{\mathfrak p\not\supseteq\mathfrak J}
      B_{\mathfrak p}
  \]
  is a finite $B$-algebra.
\end{enumerate}
Then
\[
  {\rm a)}\Longleftrightarrow{\rm b)}
  \Longrightarrow{\rm c)}
  \Longleftrightarrow{\rm d)}.
\]
Moreover, all four properties hold in either of the following
cases:
\begin{enumerate}[{\rm(i)}]
  \item $A$ is a quotient of a regular ring.
  \item $A$ is universally catenary and universally Japanese.
\end{enumerate}
\end{corollary}

\begin{proof}
Let $B=A/\mathfrak q$, where $\mathfrak q$ is a prime ideal of
$A$, so that $Y=\Spec(B)=V(\mathfrak q)$ is a closed subset of
$X$.  Put $Z=Z^{(2)}(Y)$, viewed as a subset of $X$ stable under
specialization.  Since $B$ is a domain,
$\operatorname{Ass}(A/\mathfrak q)$ consists only of the generic
point $\mathfrak q$ of $Y$.  If {\rm b)} holds, apply it to
\[
  \sh F=(A/\mathfrak q)^{\sim}
\]
and to $Z$.  Formula \sref{IV.5.9.3.1} then shows that $B^{(1)}$
is a finite $A$-module, and hence a finite $B$-module.  Thus
{\rm b)} implies {\rm a)}.

Conversely, suppose {\rm a)} holds.  Let $\sh F$ and $Z$ satisfy
the assumptions in {\rm b)}, and use the notation of
\sref{IV.5.11.1}.  Each
$Y_\alpha=\Spec(B_\alpha)$, where $B_\alpha$ is a quotient
domain of $A$.  Since
$Z_\alpha\subset Z^{(2)}(Y_\alpha)$, condition {\rm a)},
\sref{IV.5.10.2}, and the fact that
$\operatorname{Ass}(\sh O_{Y_\alpha})$ consists of the generic
point show that
\[
  \sh H^0_{Y_\alpha/Z_\alpha}(\sh O_{Y_\alpha})
\]
is coherent.  Property {\rm b)} follows from
\sref{IV.5.11.1}.

\oldpage[IV-2]{126}

To prove that {\rm c)} implies {\rm d)}, argue as above with
$\sh G=(A/\mathfrak q)^{\sim}|U$ and $T=V(\mathfrak J)$.
Conversely, {\rm d)} implies {\rm c)} by applying the same
componentwise argument and \sref{IV.5.11.1}.  Property {\rm c)}
is plainly a special case of {\rm b)}.  Finally,
\sref{IV.5.11.2}, together with the definitions of universally
catenary and universally Japanese rings, proves all four
properties under either hypothesis {\rm(i)} or {\rm(ii)}.
\end{proof}

\begin{remarks}[5.11.7]
\label{IV.5.11.7}
\begin{enumerate}[{\rm(i)}]
  \item It is not known whether condition {\rm d)} in
  \sref{IV.5.11.5} implies condition {\rm a)} without an
  additional hypothesis on $X$.  We shall see in
  \sref{IV.7.2.4} that it does when $X$ is a local scheme.
  Likewise, when $A$ is a Noetherian local ring, the four
  properties {\rm a)}, {\rm b)}, {\rm c)}, and {\rm d)} of
  \sref{IV.5.11.6} are equivalent \sref{IV.7.2.4}.  It is not
  known whether this result extends to every Noetherian ring.

  \item If $A$ satisfies property {\rm a)} of
  \sref{IV.5.11.6}, then so does every localization $S^{-1}A$
  and every finite $A$-algebra $C$.  Indeed, every quotient
  domain of $S^{-1}A$ has the form
  \[
    S^{-1}(A/\mathfrak q),
  \]
  where $\mathfrak q$ is a prime ideal of $A$ disjoint from
  $S$.  On the other hand, if $\mathfrak r$ is a prime ideal of
  $C$ and $\mathfrak p$ is its inverse image in $A$, then
  $C/\mathfrak r$ is a finite integral
  $(A/\mathfrak p)$-algebra containing $A/\mathfrak p$.  The
  assertions therefore follow from \sref{IV.5.10.17}[iii] and
  \sref{IV.5.10.17}[iv].
\end{enumerate}
\end{remarks}
